\documentclass[11pt,a4paper]{article}
\usepackage[utf8]{inputenc}
\usepackage[T1]{fontenc}
\usepackage{lmodern}
\usepackage[margin=2.2cm]{geometry}
\usepackage{graphicx}
\usepackage{booktabs}
\usepackage{url}
\usepackage[numbers,square,sort&compress]{natbib}
\usepackage{xcolor}
\usepackage{caption}
\captionsetup{font=small,labelfont=bf}
\usepackage{titling}
\setlength{\droptitle}{-2.5em}
\renewcommand{\abstractname}{\vspace{-1em}}

\title{\textbf{A profile Hidden Markov Model for the BPTI/Kunitz protease inhibitor domain built from structural multiple alignment and benchmarked against Pfam PF00014}}

\author{Amin Nami\\\small Department of Pharmacy and Biotechnology, University of Bologna, Italy}
\date{}

\begin{document}
\maketitle
\vspace{-1.5em}

\begin{abstract}
\noindent\textbf{Motivation:}
The BPTI/Kunitz domain (Pfam PF00014) is a $\sim$60-residue protease-inhibitor module of central importance in venom biology, blood coagulation and amyloid metabolism. Profile Hidden Markov Models (HMMs) are the standard tool for annotating such domains, but the impact of using a structure-based versus a sequence-based seed alignment is rarely quantified directly.

\noindent\textbf{Results:}
We built a profile HMM from a structural multiple alignment of 19 non-redundant Kunitz PDB chains and benchmarked it against the curated Pfam PF00014 HMM on the same SwissProt subset (360 positives, 574{,}267 negatives) under stratified 5-fold cross-validation. Both models achieve indistinguishable performance (MCC $0.944 \pm 0.014$ vs $0.946 \pm 0.012$) and produce identical false-positive sets, all of which are PF00014-annotated fragments excluded from the gold standard. Increasing the seed diversity to 99\% identity did not improve detection of atypical Kunitz variants, indicating that PDB coverage rather than structural alignment is the limiting factor.

\noindent\textbf{Availability:}
\url{https://github.com/AminN77/lab_bioinfo_kunitz}.

\noindent\textbf{Contact:} amin.nami@studio.unibo.it
\end{abstract}

\section{Introduction}
The BPTI/Kunitz domain is one of the most extensively characterised protease-inhibitor folds. Defined by six conserved cysteines that form three disulfide bridges, it adopts a compact $\alpha+\beta$ topology of $\sim$58 residues and inhibits serine proteases via a canonical reactive loop \citep{laskowski1980,ranasinghe2011}. The fold recurs in multi-domain blood proteins such as tissue factor pathway inhibitor (TFPI), neurotoxins (dendrotoxins), and as accessory domains in unrelated architectures including the amyloid precursor protein \citep{salvesen2013}. The corresponding Pfam family, PF00014, currently includes $\sim$10\textsuperscript{4} sequences \citep{mistry2021pfam}.

Profile Hidden Markov Models are the standard probabilistic tool for remote homology detection \citep{eddy2011}. Structural multiple alignments, when feasible, capture position-specific conservation in three-dimensional context and are expected to yield more accurate models for short, divergent domains. However, the official Pfam HMM for PF00014 is built from a curated 99-sequence seed alignment, drawing on sequence diversity that is unavailable in the PDB. In this work we ask whether a HMM built from PDB structures alone can match a sequence-curated reference, by: (i) collecting a non-redundant set of Kunitz structures from the PDB; (ii) building a profile HMM from their structural alignment with HMMER3; (iii) validating the model on SwissProt under stratified 5-fold cross-validation; and (iv) running the official Pfam HMM through the same benchmark for direct comparison.

\section{Methods}

\subsection{Structure selection and redundancy reduction}
We queried the RCSB Protein Data Bank \citep{burley2023rcsb} for polymer entities annotated with Pfam PF00014, with refinement resolution $\leq 3.0$~\AA{} and chain length 45--80 residues, yielding 238 chains. Sequences were de-duplicated with CD-HIT \citep{li2006cdhit} at 95\% identity, producing 22 representatives. Owing to the 20-structure limit of the mTM-align server, the list was trimmed to 19 chains for the main analysis. To test the effect of seed redundancy, an alternative dataset was generated by clustering at 99\% identity (58 representatives) and selecting 19 chains by greedy farthest-point sampling on pairwise sequence identity (Section~\ref{sec:control}).

\subsection{Structural alignment and HMM construction}
Single-chain PDB files (chain-specific atom records) were submitted to mTM-align \citep{dong2018mtmalign} for multiple structural alignment. The resulting FASTA alignment (57 columns for the 95\% set, 56 for the 99\%-diverse set) served as the seed for \texttt{hmmbuild} (HMMER 3.4). Self-consistency was verified by \texttt{hmmsearch} on the training sequences: all 19 entries were recovered with full-sequence E-values between $5 \times 10^{-24}$ and $2 \times 10^{-34}$.

\subsection{Pfam baseline}
The official Pfam PF00014 HMM (release Pfam 36.0, length 53 match states, 99 seed sequences, gathering threshold = 21 bits) was retrieved from InterPro \citep{paysan2023interpro}. It was searched against the same validation set with identical \texttt{hmmsearch} options (\texttt{-{}-max -E 1000}).

\subsection{Validation dataset}
The positive set comprised SwissProt entries (release 2026\_02) carrying a PF00014 cross-reference, the \texttt{reviewed} flag, and the \texttt{fragment} flag unset, queried via the UniProt REST API \citep{uniprot2023}. We removed the eight UniProt accessions associated with the 19 training PDB chains via the UniProt ID-mapping service, leaving 360 positives. The negative set consisted of the remaining 574{,}267 reviewed SwissProt entries.

\subsection{Evaluation}
Both HMMs were searched against the validation set with \texttt{hmmsearch -{}-max -E 1000}, disabling heuristic filters so that the full Forward score is computed for every sequence. Performance was assessed by stratified 5-fold cross-validation: in each fold, the optimal E-value threshold was selected on the four training folds (grid scan from $10^{-30}$ to $10^{4}$, step $10^{0.5}$) by maximising MCC, then applied to the held-out fold. Bit-score thresholds were evaluated in parallel by the same protocol. Confusion matrices, accuracy, precision, recall and MCC were computed with \texttt{scikit-learn} \citep{scikit-learn}.

\section{Results}

\subsection{Model and self-consistency}
The 95\% seed alignment spans 57 conserved match states across 19 structures and shows the canonical Kunitz signature: six conserved cysteines, a conserved Phe and Tyr in the hydrophobic core, and the GG--GG motif flanking the reactive loop. Self-search with \texttt{hmmsearch} returned all 19 training sequences with bit scores between 74.5 and 107.6 (Supplementary Table~S1).

\subsection{Performance and head-to-head comparison}
Both HMMs separate Kunitz from non-Kunitz proteins with high accuracy. The MCC--threshold curves (Fig.~\ref{fig:comp}A) are broad and flat between $10^{-12}$ and $10^{-3}$ for both models, indicating robustness to threshold choice. Precision-recall curves (Fig.~\ref{fig:comp}B) overlap almost entirely, with Pfam holding a marginal advantage at intermediate recall.

Pooled over the full benchmark at the operational threshold $E \leq 10^{-5}$, the confusion matrices are: ours --- TP 357, FP 38, FN 3, TN 574{,}229 (precision 0.904, recall 0.992, MCC 0.947); Pfam --- TP 360, FP 38, FN 0, TN 574{,}229 (precision 0.905, recall 1.000, MCC 0.951). Critically, the FP set is \emph{identical} between the two models: all 38 entries carry a PF00014 cross-reference in UniProt but were excluded from the gold standard for being fragment-flagged. Re-classifying these as positives yields a corrected MCC of 0.996 for our model and 1.000 for Pfam.

\begin{figure}[t]
\centering
\includegraphics[width=\linewidth]{comparison.pdf}
\caption{Performance comparison of our structural HMM (19 PDB chains, blue) and the official Pfam PF00014 HMM (99 seed sequences, red) on the SwissProt benchmark. (A) MCC as a function of the E-value threshold; the dashed line marks the operational threshold $10^{-5}$. (B) Precision-recall curves. The two models perform within statistical noise of each other across the full operating range.}
\label{fig:comp}
\end{figure}

\subsection{5-fold cross-validation}
Stratified 5-fold cross-validation confirms the head-to-head reading (Fig.~\ref{fig:cv}, Table~\ref{tab:cv}). Mean per-fold MCC is $0.944 \pm 0.014$ for our model and $0.946 \pm 0.012$ for Pfam, statistically indistinguishable. The standard deviation $\leq 0.015$ across folds, together with the consistency between the median optimal thresholds ($10^{-11}$ for both), indicates that neither model is overfit to a particular subset of the data.

\begin{table}[h]
\centering
\small
\caption{5-fold stratified cross-validation results. Mean $\pm$ SD across folds; the operational threshold reported is the median across folds.}
\label{tab:cv}
\begin{tabular}{lcccc}
\toprule
Model & Mean MCC & Mean Precision & Mean Recall & Median threshold \\
\midrule
Structural HMM (ours) & $0.944 \pm 0.014$ & 0.906 & 0.983 & $E \leq 10^{-11}$ \\
Pfam PF00014          & $0.946 \pm 0.012$ & 0.909 & 0.986 & $E \leq 10^{-10}$ \\
\bottomrule
\end{tabular}
\end{table}

\begin{figure}[h]
\centering
\includegraphics[width=0.75\linewidth]{cv5.pdf}
\caption{Per-fold MCC across the 5-fold stratified cross-validation. Dashed lines indicate the mean over folds.}
\label{fig:cv}
\end{figure}

\subsection{Bit-score versus E-value thresholds}
The Pfam HMM is published with a curated gathering threshold of 21 bits, independent of database size. We re-ran the cross-validation using bit-score thresholds. Every fold converged to a bit-score threshold of $\approx 21$ for the Pfam model (SD 0.9 bits), recovering recall = 0.997 at MCC = 0.949. For our model the optimal bit-score threshold was more variable ($21 \pm 6$ bits) and yielded a slightly lower mean MCC of 0.938. The greater stability of the Pfam bit-score reflects the larger and more diverse seed alignment used by Pfam curators. Because bit scores are independent of database size, they are the recommended portable threshold when the search target may grow, e.g. when transferring from SwissProt to TrEMBL.

\subsection{False positive and false negative analysis}
All 38 false positives at $E \leq 10^{-5}$ are SwissProt entries carrying a PF00014 cross-reference, including TFPI (P10646), the amyloid precursor protein (P05067), inter-$\alpha$-trypsin inhibitor (P04365), AMBP (P02760), collagen $\alpha$3(VI) (P12111), and a panel of snake/spider/anemone toxin Kunitz-like inhibitors. Their exclusion from the positive set traces to the SwissProt \texttt{fragment} flag (typically venom peptides obtained by direct protein sequencing of incomplete chains). Both our HMM and the Pfam HMM flag the \emph{same} 38 entries, providing external evidence that they are true Kunitz domains rather than spurious hits.

Three positives fall below the threshold for our model: D3GGZ8 and O62247 (\emph{C.\ elegans} Kunitz-type protein bli-5 and its paralog, with non-canonical cysteine spacing and an additional non-Kunitz domain in O62247) and A0A1Q1NL17 (snake anticoagulant HA11, with substituted reactive-loop residues). Pfam detects all three at $E \approx 10^{-7}$, suggesting that the breadth of the sequence-curated Pfam seed is the principal source of the residual gap.

\subsection{Effect of seed redundancy --- a control experiment}
\label{sec:control}
We tested whether the residual gap could be closed by using a more diverse structural seed. CD-HIT at 99\% identity produced 58 PDB chains; greedy farthest-point sampling on sequence identity selected 19 maximally diverse representatives (adding four PDB chains not used in the main model: 6HAR\_E, 6BX8\_D, 1YLD\_B, 4ISN\_B). The resulting HMM (\textsc{V2}) achieved an essentially identical pooled MCC of 0.945 versus 0.947 for the 95\% seed (\textsc{V1}) and did \emph{not} recover the three atypical FNs. The PDB chain-length filter (45--80 residues) selects against precisely the divergent Kunitz variants present in venoms and multi-domain proteins, which are over-represented in the Pfam seed. This indicates that, for this family, PDB coverage is the limiting factor for structural-MSA-based HMMs, not the diversity of the structural alignment itself.

\section{Discussion}
A profile HMM trained on a 19-chain structural alignment of Kunitz domains achieves performance statistically indistinguishable from the official Pfam HMM built from a 99-sequence curated alignment. The two HMMs make identical errors on the SwissProt benchmark, and the residual three false negatives correspond to divergent venom and nematode Kunitz variants whose sequence variation lies outside the PDB-accessible chain-length window. Increasing structural diversity through a 99\%-identity CD-HIT cut does not improve detection of these variants, suggesting that PDB coverage rather than structural-alignment quality is the limiting factor. This is a structural argument for combining sequence and structure information --- as Pfam in fact does in its seed curation --- rather than relying on either source alone for short, taxonomically dispersed domain families.

The MCC--threshold plateau spanning nine orders of magnitude in E-value, and the agreement between independent 5-fold CV partitions (SD $< 0.015$), indicate that the model is well calibrated and not overfit. The bit-score analysis underscores that, while E-value thresholds are convenient and database-specific, the Pfam-style curated bit-score threshold ($\approx 21$ bits) is more stable across folds and is recommended for cross-database transfer.

\section{Conclusion}
We provide a compact, accurate Kunitz HMM (57 match states, 19 PDB seed chains) that matches the curated Pfam reference within statistical noise, together with a complete validation pipeline. The model is suitable as a drop-in tool for proteome-wide Kunitz annotation, and the comparison framework can be reused to benchmark any custom HMM against the corresponding Pfam reference under identical conditions.

\section*{Availability and Implementation}
Source code, seed alignments, trained HMMs, the Pfam baseline, and the evaluation pipeline are available at \url{https://github.com/AminN77/lab_bioinfo_kunitz}. The pipeline runs on Linux with HMMER~3.4, CD-HIT~4.8.1, Python~3.12, Biopython~1.83 and scikit-learn~1.4.

\bibliographystyle{unsrtnat}
\begin{thebibliography}{99}

\bibitem{eddy2011} Eddy, S.~R. (2011). Accelerated profile HMM searches. \textit{PLoS Comput.\ Biol.}, 7(10), e1002195.

\bibitem{dong2018mtmalign} Dong, R., Pan, S., Peng, Z., Zhang, Y., Yang, J. (2018). mTM-align: a server for fast protein structure database search and multiple protein structure alignment. \textit{Nucleic Acids Res.}, 46(W1), W380--W386.

\bibitem{burley2023rcsb} Burley, S.~K. \emph{et al.} (2023). RCSB Protein Data Bank (RCSB.org): delivery of experimentally-determined PDB structures alongside one million computed structure models of proteins from artificial intelligence/machine learning. \textit{Nucleic Acids Res.}, 51(D1), D488--D508.

\bibitem{mistry2021pfam} Mistry, J. \emph{et al.} (2021). Pfam: The protein families database in 2021. \textit{Nucleic Acids Res.}, 49(D1), D412--D419.

\bibitem{paysan2023interpro} Paysan-Lafosse, T.\ \emph{et al.} (2023). InterPro in 2022. \textit{Nucleic Acids Res.}, 51(D1), D418--D427.

\bibitem{uniprot2023} The UniProt Consortium (2023). UniProt: the Universal Protein Knowledgebase in 2023. \textit{Nucleic Acids Res.}, 51(D1), D523--D531.

\bibitem{li2006cdhit} Li, W., Godzik, A. (2006). Cd-hit: a fast program for clustering and comparing large sets of protein or nucleotide sequences. \textit{Bioinformatics}, 22(13), 1658--1659.

\bibitem{laskowski1980} Laskowski, M., Kato, I. (1980). Protein inhibitors of proteinases. \textit{Annu.\ Rev.\ Biochem.}, 49, 593--626.

\bibitem{ranasinghe2011} Ranasinghe, S., McManus, D.~P. (2013). Structure and function of invertebrate Kunitz serine protease inhibitors. \textit{Dev.\ Comp.\ Immunol.}, 39(3), 219--227.

\bibitem{salvesen2013} Rawlings, N.~D., Barrett, A.~J., Bateman, A. (2010). MEROPS: the peptidase database. \textit{Nucleic Acids Res.}, 38(Database issue), D227--D233.

\bibitem{scikit-learn} Pedregosa, F. \emph{et al.} (2011). Scikit-learn: Machine Learning in Python. \textit{J.\ Mach.\ Learn.\ Res.}, 12, 2825--2830.

\end{thebibliography}

\end{document}
